\documentclass[12pt]{article}
\usepackage{amsmath,amssymb,amsfonts}
\usepackage[margin=1in]{geometry}

\begin{document}

\begin{center}
{\Large\bfseries Chapter 10, Problem 10}
\end{center}

\noindent\textbf{Problem.}
\begin{itemize}
\item[a)] By direct calculation, show that
\[
\nabla^2 = \frac{\partial^2}{\partial x^2} + \frac{\partial^2}{\partial y^2} + \frac{\partial^2}{\partial z^2}
\]
is invariant under rotation through the Euler angles $\alpha$, $\beta$ and $\gamma$. To refresh your memory about Euler angles, see Section 1.4. [Hint: Note that the most general rotation is built up from three simpler rotations.]

\item[b)] Determine how the three components of the quantum mechanical angular momentum [see Eqs.~(5.86)] transform under the group of rotations about the $z$-axis.
\end{itemize}

\bigskip
\noindent\textbf{Solution.}

\noindent\textbf{Part (a):}

A general rotation specified by Euler angles $(\alpha, \beta, \gamma)$ is built up from three successive rotations:
\begin{enumerate}
\item Rotation by $\alpha$ about the $z$-axis
\item Rotation by $\beta$ about the new $y'$-axis (or the line of nodes)
\item Rotation by $\gamma$ about the new $z''$-axis
\end{enumerate}

It suffices to show $\nabla^2$ is invariant under each of these. By symmetry it is enough to show invariance under a rotation about one axis, say the $z$-axis by angle $\theta$.

Under rotation about the $z$-axis by angle $\theta$:
\[
x' = x\cos\theta + y\sin\theta, \quad y' = -x\sin\theta + y\cos\theta, \quad z' = z.
\]

By the chain rule:
\[
\frac{\partial}{\partial x} = \cos\theta\,\frac{\partial}{\partial x'} - \sin\theta\,\frac{\partial}{\partial y'}, \quad
\frac{\partial}{\partial y} = \sin\theta\,\frac{\partial}{\partial x'} + \cos\theta\,\frac{\partial}{\partial y'}, \quad
\frac{\partial}{\partial z} = \frac{\partial}{\partial z'}.
\]

Computing:
\begin{align*}
\frac{\partial^2}{\partial x^2} &= \cos^2\theta\,\frac{\partial^2}{\partial x'^2} - 2\sin\theta\cos\theta\,\frac{\partial^2}{\partial x'\partial y'} + \sin^2\theta\,\frac{\partial^2}{\partial y'^2}, \\[6pt]
\frac{\partial^2}{\partial y^2} &= \sin^2\theta\,\frac{\partial^2}{\partial x'^2} + 2\sin\theta\cos\theta\,\frac{\partial^2}{\partial x'\partial y'} + \cos^2\theta\,\frac{\partial^2}{\partial y'^2}.
\end{align*}

Adding:
\[
\frac{\partial^2}{\partial x^2} + \frac{\partial^2}{\partial y^2} = (\cos^2\theta + \sin^2\theta)\frac{\partial^2}{\partial x'^2} + (\sin^2\theta + \cos^2\theta)\frac{\partial^2}{\partial y'^2} = \frac{\partial^2}{\partial x'^2} + \frac{\partial^2}{\partial y'^2}.
\]

Therefore:
\[
\nabla^2 = \frac{\partial^2}{\partial x^2} + \frac{\partial^2}{\partial y^2} + \frac{\partial^2}{\partial z^2} = \frac{\partial^2}{\partial x'^2} + \frac{\partial^2}{\partial y'^2} + \frac{\partial^2}{\partial z'^2} = \nabla'^2.
\]

The same argument applies to rotations about the $y$-axis (with $x$ and $z$ mixing while $y$ is fixed). Since the general rotation is a product of three such rotations, and the Laplacian is invariant under each, it is invariant under the general rotation. $\square$

\bigskip
\noindent\textbf{Part (b):}

The quantum mechanical angular momentum components are:
\[
L_x = -i\hbar\!\left(y\frac{\partial}{\partial z} - z\frac{\partial}{\partial y}\right), \quad
L_y = -i\hbar\!\left(z\frac{\partial}{\partial x} - x\frac{\partial}{\partial z}\right), \quad
L_z = -i\hbar\!\left(x\frac{\partial}{\partial y} - y\frac{\partial}{\partial x}\right).
\]

Under a rotation by angle $\theta$ about the $z$-axis:
\[
x' = x\cos\theta + y\sin\theta, \quad y' = -x\sin\theta + y\cos\theta, \quad z' = z.
\]

Expressing in primed coordinates:
\begin{align*}
x &= x'\cos\theta - y'\sin\theta, \\
y &= x'\sin\theta + y'\cos\theta, \\
z &= z'.
\end{align*}

And the partial derivatives transform as:
\begin{align*}
\frac{\partial}{\partial x'} &= \cos\theta\,\frac{\partial}{\partial x} + \sin\theta\,\frac{\partial}{\partial y}, \\
\frac{\partial}{\partial y'} &= -\sin\theta\,\frac{\partial}{\partial x} + \cos\theta\,\frac{\partial}{\partial y}, \\
\frac{\partial}{\partial z'} &= \frac{\partial}{\partial z}.
\end{align*}

Now compute how $L_x$ and $L_y$ transform. In the rotated frame:
\begin{align*}
L_x' &= -i\hbar\!\left(y'\frac{\partial}{\partial z'} - z'\frac{\partial}{\partial y'}\right), \\
L_y' &= -i\hbar\!\left(z'\frac{\partial}{\partial x'} - x'\frac{\partial}{\partial z'}\right), \\
L_z' &= -i\hbar\!\left(x'\frac{\partial}{\partial y'} - y'\frac{\partial}{\partial x'}\right).
\end{align*}

Substituting back to the unprimed variables:
\begin{align*}
L_x' &= \cos\theta\, L_x + \sin\theta\, L_y, \\
L_y' &= -\sin\theta\, L_x + \cos\theta\, L_y, \\
L_z' &= L_z.
\end{align*}

In matrix form:
\[
\begin{pmatrix} L_x' \\ L_y' \\ L_z' \end{pmatrix} =
\begin{pmatrix} \cos\theta & \sin\theta & 0 \\ -\sin\theta & \cos\theta & 0 \\ 0 & 0 & 1 \end{pmatrix}
\begin{pmatrix} L_x \\ L_y \\ L_z \end{pmatrix}.
\]

The angular momentum components \textbf{transform as a vector} (i.e., they transform under the same rotation matrix that acts on the coordinates). In particular, $L_z$ is invariant under rotations about the $z$-axis, while $L_x$ and $L_y$ mix among themselves.

\end{document}
